\documentclass{article}
\usepackage[utf8]{inputenc}
\usepackage[T1]{fontenc}
\usepackage{geometry}
\usepackage{graphicx}
\usepackage{float}
\usepackage{hyperref}
\usepackage[english,indonesian]{babel}

\title{Optimasi Pemrosesan Template LaTeX Menggunakan Pendekatan Berbasis AI}
\author{Academic Support Agent}
\date{\today}

\begin{document}
\maketitle

\begin{abstract}

\end{abstract}

\chapter{PENDAHULUAN}


\chapter{TINJAUAN PUSTAKA}


% Original content (commented out):
% \documentclass[10pt]{article}
\usepackage[utf8]{inputenc}
\usepackage[margin=5mm]{geometry}

\include{planner-template}

\title{100 Days Challenge Template}
\author{The \LaTeX{} Ninja}
\date{October 2020}



\begin{document}
%--------------------------------------------------------------------------------------------
% use out of the box
%--------------------------------------------------------------------------------------------    
\onehundreddaychallengetracker{\#100DaysofCode}{\faLaptop}

%--------------------------------------------------------------------------------------------
% adapt to a 30 day challenge
%--------------------------------------------------------------------------------------------    
\newpage 

\section{Health-Challenge \hfill Start: \hspace{3cm}}

\smallsymboltracker{Health
}{\faLeaf}{greyedout}{30}{0.7} % Symbol - Colour - Number of Bubbles - percentage of textwidth to take up
\vspace{3em}

\plannerheading{Rewards for completion}

\vspace{5em}
\plannerheading{Big project for the whole challenge}

\vspace{5em}
\plannerheading{Subprojects}

\newpage
\tendayschallenge{10}{Yoga}{\faChild}
\tendayschallenge{10}{Eating Well}{\faCutlery}
\tendayschallenge{10}{Meditation}{\faLeaf}



\end{document}


\end{document}